\documentclass[11pt]{article}
\usepackage[margin=1in]{geometry}
\usepackage{amsmath,amssymb}
\usepackage{mathtools}
\usepackage{enumitem}
\usepackage{hyperref}
\usepackage{titlesec}
\usepackage{parskip}

\title{Reinforcement Learning from Scratch:\\
Policy Gradients, REINFORCE, Actor-Critic, and PPO\\
\large A primer for \texttt{rl\_train.py}}
\author{}
\date{}

\begin{document}
\maketitle

\section*{The setup: what problem RL is even solving}

Everything below assumes an \textbf{MDP (Markov Decision Process)} --- the standard
mathematical model for sequential decision-making.

\paragraph{Agent and environment.} There's an \emph{agent} (the thing being trained ---
here, the policy network) and an \emph{environment} (the CUDA particle sim). They interact
in discrete time steps $t = 0, 1, 2, \dots$

\begin{itemize}
  \item $S_t$ --- the \textbf{state} at time $t$: whatever information the agent observes
  about the environment. In this codebase, $S_t$ is the occupancy grid.
  \item $A_t$ --- the \textbf{action} the agent takes at time $t$, chosen after observing
  $S_t$. Here, $A_t$ is the force-grid field.
  \item $R_{t+1}$ --- the \textbf{reward}: a scalar the environment returns after the agent
  takes action $A_t$ in state $S_t$, indicating how good that was. Note the indexing
  convention: $R_{t+1}$ is the consequence of $(S_t, A_t)$, not $R_t$ --- the reward arrives
  one tick after the state/action that caused it. Here,
  $R_{t+1} = -\bar d - 0.1\,\sigma_d$ (the distance-to-center penalty).
  \item The environment then produces $S_{t+1}$, and the cycle repeats.
\end{itemize}

This handoff --- $S_t, A_t \to R_{t+1}, S_{t+1}$ --- is a \textbf{transition}. A full run
from start to termination, $S_0, A_0, R_1, S_1, A_1, R_2, \dots$, is an \textbf{episode} or
\textbf{trajectory}, often written $\tau$.

\textbf{``Markov''} means $S_{t+1}$ and $R_{t+1}$ depend only on $(S_t, A_t)$, not on
anything earlier. This is a modeling assumption, not automatically true: an occupancy-only
observation (no velocity) technically violates it, since two states with identical
positions but different velocities look identical to the agent but evolve differently.

The \textbf{dynamics} are written
$$p(s', r \mid s, a) \doteq \Pr\{S_{t+1}=s', R_{t+1}=r \mid S_t=s, A_t=a\},$$
the probability of landing in state $s'$ with reward $r$, given the agent was in state $s$
and took action $a$. This belongs to the environment; the agent never sees it directly ---
in this case it's the CUDA kernels, and the agent only observes samples of the outcome.

\section*{The policy}

$\pi(a \mid s)$ --- the \textbf{policy}: a conditional probability distribution over
actions given a state, $\pi(a|s) = \Pr\{A_t = a \mid S_t = s\}$. This is the object being
learned. When implemented by a parameterized function (a neural net) with parameter vector
$\theta$, write $\pi(a|s,\theta)$ to make the dependence on the learnable parameters
explicit --- $\theta$ is just ``all the weights,'' stacked into one vector. In the code,
$\theta$ is every weight in \texttt{ActorCritic}'s conv layers plus \texttt{log\_std}.

A policy can be \textbf{deterministic} (one action per state, probability 1) or
\textbf{stochastic} (a real distribution). Differentiating through a distribution's
\emph{parameters} (mean/std of a Gaussian, or softmax logits) is what makes gradient-based
learning of $\pi$ possible at all; a hard deterministic policy $a=f(s)$ has no
``probability of the action taken'' to differentiate. (Deterministic policy gradients
exist too --- Silver et al.\ 2014 --- but need a different trick, differentiating through
a learned $Q$ function instead; not covered here.)

\section*{Return: ``good'' over time, not just at one step}

A single reward $R_{t+1}$ only tells you about one step. Define the \textbf{return}:
$$G_t \doteq R_{t+1} + \gamma R_{t+2} + \gamma^2 R_{t+3} + \cdots = \sum_{k=0}^\infty \gamma^k R_{t+k+1}.$$
$\gamma \in [0,1]$ is the \textbf{discount factor} --- how much a reward $k$ steps in the
future is worth relative to right now. $\gamma$ near 1: farsighted. $\gamma=0$: only the
immediate reward matters. It guarantees $G_t$ is finite even for long/continuing tasks,
and is often a genuinely correct model of preference. In \texttt{rl\_train.py},
\texttt{gamma=0.99} in \texttt{compute\_gae}.

$G_t$ is a \textbf{random variable} --- it depends on the random future trajectory --- so
any statement about it is really a statement about its expectation or distribution.

\section*{Value functions}

\begin{align*}
v_\pi(s) &\doteq \mathbb E_\pi[G_t \mid S_t = s] &&\text{\textbf{state-value function}: expected return from state $s$ under $\pi$}\\
q_\pi(s,a) &\doteq \mathbb E_\pi[G_t \mid S_t = s, A_t = a] &&\text{\textbf{action-value function} (``Q-function''): expected return, taking}\\
&&&\text{action $a$ in $s$ \emph{right now}, then following $\pi$ afterward}
\end{align*}
$\mathbb E_\pi[\cdot]$ means expectation where actions are drawn from $\pi$ and states
evolve under the environment's real dynamics.

Both are unknown functions --- defined in terms of the future, which hasn't happened yet.
Estimating them is the core of \emph{action-value methods}: learn $\hat q(s,a)$, then act
by $\arg\max_a \hat q(s,a)$, so the policy is \emph{derived from} the value function.
Policy-gradient methods break from this: they parameterize and learn $\pi$ directly, using
a value function only as a helper signal (the ``critic''), never as what action-selection
reads from.

$\hat v(s,w)$, $\hat q(s,a,w)$ --- a \textbf{learned approximation} to $v_\pi$ or $q_\pi$,
with its own parameter vector $w$ (distinct from the policy's $\theta$) --- e.g.\ the critic
head of \texttt{ActorCritic}.

\section*{The learning objective and gradient ascent}

$$J(\theta) \doteq v_{\pi_\theta}(s_0)$$
the \textbf{performance measure}: for the episodic case, the value of a fixed start state
$s_0$ under the policy that $\theta$ induces. This is the number being maximized by
choosing $\theta$. (A continuing-task version, $J(\theta) = r(\pi)$, the average reward
rate, also exists --- relevant if the task never terminates. Episodes here terminate on
domain-exit, so the episodic version applies, modulo the finite-horizon truncation from
\texttt{episode\_steps}.)

\textbf{Gradient ascent}: repeatedly nudge $\theta$ in the direction that increases $J$
fastest:
$$\theta_{t+1} = \theta_t + \alpha \nabla J(\theta_t).$$
$\alpha$ is the \textbf{step size} (learning rate). $\nabla J(\theta)$ is the
\textbf{gradient} of $J$ w.r.t.\ $\theta$ --- a vector of partial derivatives,
$\nabla J(\theta) = \big[\partial J/\partial\theta_1,\ \partial J/\partial\theta_2,\dots\big]^\top$,
pointing toward steepest increase.

\textbf{Problem}: $\nabla J(\theta)$ can't be computed exactly --- that would require
knowing the environment's dynamics $p(s',r|s,a)$ and summing over every possible
trajectory. You only get to \emph{sample} trajectories by running the policy. So instead
compute $\widehat{\nabla J(\theta_t)}$, a \textbf{stochastic estimate} --- something
calculable from one sampled trajectory, whose expected value equals (or is proportional
to) the true gradient:
$$\mathbb E\big[\widehat{\nabla J(\theta_t)}\big] \propto \nabla J(\theta_t).$$
This is \textbf{stochastic gradient ascent} --- same idea as SGD in supervised learning,
applied to a much less well-behaved objective: there's no labeled data, the agent
generates its own data by acting, and that data's distribution depends on the very
parameters being updated. Everything that follows is: \emph{what is a valid, computable}
$\widehat{\nabla J(\theta_t)}$?

\section*{The Policy Gradient Theorem, symbol by symbol}

$$\nabla J(\theta) \propto \sum_s \mu(s) \sum_a q_\pi(s,a)\,\nabla \pi(a|s,\theta)$$

\begin{itemize}
  \item $\sum_s(\cdot)$, $\sum_a(\cdot)$ --- sums over every possible state / action (for
  continuous spaces these become integrals; the continuous-action Gaussian policy below is
  the relevant generalization).
  \item $\mu(s)$ --- the \textbf{on-policy state distribution}: roughly, the fraction of
  time steps the agent spends in state $s$ if it runs $\pi$ for a long time. It's an
  emergent property of $\pi$ and the environment, not something chosen directly.
  \item $q_\pi(s,a)$ --- as above.
  \item $\nabla\pi(a|s,\theta)$ --- gradient of the policy's own output probability w.r.t.\
  $\theta$. This \emph{is} directly computable once a parameterization is chosen (plain
  calculus on the network/softmax/Gaussian formula), unlike $\nabla J$ itself.
  \item $\propto$ --- ``proportional to'': the two sides differ by a $\theta$-independent
  constant (for episodic tasks, the average episode length), so it doesn't matter for
  gradient \emph{ascent} --- any constant can be absorbed into $\alpha$.
\end{itemize}

\textbf{Why this theorem is the crux of the chapter}: naively, $J(\theta)$ depends on
$\theta$ through two channels --- (1) $\theta$ changes which actions get taken in a given
state (easy: that's just $\nabla\pi$), and (2) $\theta$ changes \emph{which states the
agent tends to visit at all}, and that channel's dependence on $\theta$ runs through the
environment's unknown dynamics. The theorem's content is that the formula for
$\nabla J(\theta)$ contains \textbf{no derivative of $\mu$} --- channel (2)'s messiness
cancels out algebraically (via an infinite-unrolling argument). That's what makes the
gradient computable from experience alone, without a model of the environment.

\section*{From theorem to something runnable: REINFORCE}

The theorem is a sum weighted by $\mu(s)$ --- exactly ``how often you visit $s$ running
$\pi$.'' So instead of computing the sum literally, \textbf{sample} it: run the policy, and
the states/actions encountered already occur in proportion to $\mu(s)\pi(a|s)$. Multiplying
and dividing the inner sum by $\pi$ turns it into an expectation over the actually-sampled
action $A_t$; using $\mathbb E_\pi[G_t \mid S_t,A_t] = q_\pi(S_t,A_t)$ (the definition of
$q_\pi$) gives:
$$\nabla J(\theta) \propto \mathbb E_\pi\left[G_t\,\frac{\nabla\pi(A_t|S_t,\theta)}{\pi(A_t|S_t,\theta)}\right].$$
Everything inside is computable from one played-out episode: $G_t$ is the observed
discounted sum of future rewards, and
$\nabla\pi(A_t|S_t,\theta)/\pi(A_t|S_t,\theta) = \nabla\ln\pi(A_t|S_t,\theta)$
(since $\nabla\ln x = \nabla x / x$) is computable by autodiff given the policy's formula.
This vector is the \textbf{eligibility vector} --- the direction in $\theta$-space that
most increases the probability of the action just taken. The \textbf{REINFORCE} update:
$$\theta_{t+1} = \theta_t + \alpha\, G_t\, \nabla\ln\pi(A_t|S_t,\theta_t).$$
In words: push $\theta$ toward making $A_t$ more likely in state $S_t$, scaled by how good
the return turned out to be. Good outcome $\Rightarrow$ reinforce that action; bad
(very negative) outcome $\Rightarrow$ push away from it.

\textbf{``Monte Carlo''} here means $G_t$ is computed by summing real observed future
rewards to the end of the episode, rather than estimated/bootstrapped from a model. This
makes REINFORCE unbiased but high-variance (one lucky or unlucky episode swings $G_t$
wildly) and usable only once a full episode has finished.

\section*{Baselines: cutting variance without changing the target}

Any function of state alone, $b(s)$, can be subtracted from $q_\pi(s,a)$ inside the
theorem without changing its value:
$$\nabla J(\theta) \propto \sum_s\mu(s)\sum_a\big(q_\pi(s,a)-b(s)\big)\nabla\pi(a|s,\theta).$$
This holds because
$$\sum_a b(s)\nabla\pi(a|s,\theta) = b(s)\,\nabla\!\left(\sum_a\pi(a|s,\theta)\right) = b(s)\,\nabla(1) = 0:$$
probabilities over actions in a state must sum to 1 regardless of $\theta$, so its gradient
is the zero vector --- multiplying by anything (here $b(s)$) and adding it in contributes
exactly nothing, in expectation. The natural choice is $b(s) = \hat v(s,w)$ --- a learned
estimate of how good state $s$ is in general. Then $G_t - \hat v(S_t,w)$ measures ``how
much better/worse did this episode do than expected from this state'' --- much lower
variance than $G_t$ alone, since it's \emph{relative} to a per-state expectation rather
than an absolute number.

\section*{Actor-critic: bootstrapping instead of waiting for the full return}

REINFORCE-with-baseline still needs the full $G_t$ (wait for episode end).
\textbf{Bootstrapping} means: instead of summing real future rewards to termination, use
the current value estimate as a stand-in for ``everything after the next step'':
$$G_{t:t+1} \doteq R_{t+1} + \gamma\hat v(S_{t+1}, w),$$
the \textbf{one-step return} or \textbf{TD (temporal-difference) target}. Define the
\textbf{TD error}:
$$\delta_t \doteq R_{t+1} + \gamma\hat v(S_{t+1},w) - \hat v(S_t,w),$$
literally ``reward received, plus a discounted estimate of what's still to come, minus
what this state was previously thought to be worth.'' Substituting $\delta_t$ for
$G_t - \hat v(S_t,w)$:
$$\theta_{t+1} = \theta_t + \alpha\,\delta_t\,\nabla\ln\pi(A_t|S_t,\theta_t).$$
This is fully \textbf{online} (updates every step, no waiting for episode end) --- the
tradeoff is that $\delta_t$ is \emph{biased} (it depends on $\hat v$, itself imperfect and
still being learned), but far lower variance than the Monte Carlo $G_t$. When the value
function is used this way --- to critique the action just taken, rather than merely as a
baseline for an already-fully-observed return --- it's the \textbf{critic}, and $\pi_\theta$
is the \textbf{actor}; the combination is an \textbf{actor-critic method}. This
bias/variance dial (one-step TD $\leftrightarrow$ full Monte Carlo) is exactly what GAE's
$\lambda$ interpolates continuously --- GAE is the same $\delta_t$, generalized.

\textbf{Eligibility traces} (using accumulators $z_\theta, z_w$): a running
exponentially-decayed accumulator of past gradients/features, letting a TD-style method
credit not just the immediately-preceding state/action but a decaying window of recent
ones --- interpolating smoothly between the one-step update above and the full-episode
Monte Carlo update, without storing the whole episode. Not used in \texttt{rl\_train.py}
(GAE serves a related but distinct purpose there, computed after the fact over a stored
rollout buffer rather than online via a decaying trace).

\section*{Continuous actions --- matches the Gaussian policy in code}

Everything above assumed a \emph{discrete} action, so $\pi(a|s,\theta)$ was literally a
probability. For a continuous (real-valued) action, $\pi(a|s,\theta)$ is instead a
\textbf{probability density function} --- $\pi(a|s,\theta)\,da$ gives the probability the
action lands in a small interval of width $da$ around $a$; $\pi$ itself can exceed 1 (a
density, not a probability). The canonical example is a \textbf{Gaussian/Normal policy}:
$$\pi(a|s,\theta) = \frac{1}{\sigma(s,\theta)\sqrt{2\pi}}\exp\left(-\frac{(a-\mu(s,\theta))^2}{2\sigma(s,\theta)^2}\right).$$
$\mu(s,\theta)$, $\sigma(s,\theta)$ --- the \textbf{mean} and \textbf{standard deviation} of
the action distribution, each a parameterized function of state (simplest case: linear in
features; in this codebase, a CNN for $\mu$ and a free, state-independent parameter for
$\log\sigma$). Everything from REINFORCE through actor-critic goes through unchanged ---
just plug this $\pi$'s eligibility vector, $\nabla\ln\pi(a|s,\theta)$, into the same update
rules. This is precisely what \texttt{torch.distributions.Normal} plus
\texttt{.log\_prob().sum()} implements: PyTorch computes $\nabla\ln\pi(a|s,\theta)$ via
autodiff instead of deriving the closed form by hand.

\section*{Two loose ends needed for PPO specifically}

\begin{itemize}
  \item \textbf{Advantage}: $A_\pi(s,a) \doteq q_\pi(s,a) - v_\pi(s)$ --- the
  baseline-subtracted quantity from the baseline section above, given its own name and
  symbol because it becomes the central object in later work like GAE/PPO.
  \item \textbf{Importance sampling ratio}:
  $$r_t(\theta) \doteq \frac{\pi_\theta(A_t|S_t)}{\pi_{\theta_{\text{old}}}(A_t|S_t)},$$
  a correction factor letting data collected under an \emph{old} policy $\theta_{\text{old}}$
  be reused to estimate a gradient for the \emph{current} $\theta$, by reweighting each
  sample by how much more (or less) likely the new policy is to have taken that same
  action. This is what lets PPO run several epochs of gradient steps on one batch of
  rollout data instead of discarding it after a single REINFORCE-style update.
\end{itemize}

\end{document}
